\documentclass[a4]{article}

% PAQUETES %%%%%%%%%%%%%%%%%%%%%%%%%%%%%%%%%%%%%%%%%%%%%%%%%%%%%%%%%%%%%%%%%%%%%
\usepackage[utf8]{inputenc}
\usepackage{framed} % Para crear cuadros alrededor del texto
\usepackage{makecell}
\usepackage[english]{babel}
%
\usepackage[fixlanguage]{babelbib}
    \bibliographystyle{IEEEtran}

\usepackage{url}

%!TEX encoding = UTF-8 Unicode
% !TEX spellcheck = es

\setlength{\textwidth}{6.5 in}
\setlength{\textheight}{8.5 in}
\setlength{\headsep}{0.5 in}

\usepackage{multirow, array} % para las tablas
\usepackage{float} % para usar [H]
\usepackage{charter}

\usepackage[breaklinks=true]{hyperref}

%
% ADD PACKAGES here:
%

\usepackage[utf8]{inputenc}
\usepackage{amsmath,longtable,siunitx, float, graphicx, dcolumn, bm, fancyhdr, authblk, subcaption, caption, hyperref, multicol, setspace}

\usepackage{bbding}
\usepackage{amssymb}
\usepackage[rightcaption]{sidecap}
%\usepackage{textcomp}

\usepackage[hmarginratio=1:1, top=32mm, columnsep=20pt, headsep=\dimexpr20mm, margin=0.7in, top=1in ]{geometry}

\usepackage{booktabs}
\usepackage{multirow}

\usepackage{algorithm}
\usepackage{enumitem}
\usepackage{tabularx}
\usepackage[noend]{algpseudocode}
\usepackage{xcolor}

\usepackage{tikz}
\usetikzlibrary{shapes.geometric, arrows.meta, positioning}

% ----------- global styles -----------
\tikzset{
  font=\scriptsize,                      % tiny text everywhere
  every node/.style       = {transform shape},   % respect scale
  startstop/.style        = {ellipse, draw, fill=blue!10,
                             inner sep=1pt, minimum width=13mm},
  process/.style          = {rectangle, draw, fill=orange!15,
                             text width=26mm, inner sep=2pt, align=center},
  decision/.style         = {diamond, draw, fill=green!15, aspect=2.3,
                             text width=20mm, inner sep=1pt, align=center},
  arrow/.style            = {->, >=Stealth, shorten >=1pt, shorten <=1pt}
}

\newcommand{\TODO}{\textbf{\textcolor{red}{TODO}}}

\renewcommand{\algorithmicrequire}{\textbf{Input:}}
\renewcommand{\algorithmicensure}{\textbf{Output:}}

\lhead{\includegraphics[scale=0.4]{CETC.png}}
\rhead{\includegraphics[scale=0.025]{Cornell-University-Logo.png}}

\newtheorem{theo}{Theorem}
\newenvironment{theorem}
  {\begin{mdframed}\begin{theo}}
  {\end{theo}\end{mdframed}}
  
\newtheorem{prop}{Proposición}
\newenvironment{proposition}
  {\begin{mdframed}\begin{prop}}
  {\end{prop}\end{mdframed}}

\newtheorem{lem}{Lema}
\newenvironment{lema}
  {\begin{mdframed}\begin{lem}}
  {\end{lem}\end{mdframed}}
  
\newtheorem{defi}{Definition}
\newenvironment{definition}
  {\begin{mdframed}\begin{def}}
  {\end{def}\end{mdframed}}
  
\newtheorem{proof}{Proof}
\newtheorem{remark}{Remark}

\usepackage[numbers]{natbib}
\renewcommand{\thesection}{\Roman{section}}
%\setlength{\parskip}{\baselineskip}%

\begin{document}

\thispagestyle{fancy}

\begin{center}
    \textbf{\large Orbit Odyssey - Computer Vision} \\
    {Scope of Work - Updated Phase 2 and 3}
    
    \vspace{0.7em}

    David Alejandro Fuquen \\

    {\small \today}

\end{center}

\section*{Scope of Work for the Orbit Odyssey - Computer Vision Autonomous Challenge - Phase 2 and Phase 3}

\subsection*{Summary}
Updated scope of works for Phase 2 and Phase 3, updated after the completion of Phase 1.



%===============================================================================
\section*{Phase 2: Graphical User Interface}
%===============================================================================

\noindent \textbf{Estimated hours:} 42 \hfill 

\vspace{0.5em}

This phase creates a new, web-based GUI to work with CETC's current codebase. It also implements a cryptographic anti-cheating system: a signed, encrypted receipt that records the exact hyperparameters and results from each training run. Only moderators possess the decryption key.

\subsection*{Scope of Work}

\begin{enumerate}[leftmargin=*]
    \item \textbf{Planning and documentation} (5 hrs)
    \begin{itemize}
        \item Define the 3 object used for the Computer Vision Challenege, e.g. which 3D print to use, what size, and which tin. These must be easily available worldwide.
        \item Adapt and complete the Computer Vision presentation (PPT 7.X within the google drive folder)
    \end{itemize}
    \item \textbf{CETC Code Integration \& GUI creation} (15 hrs)
    \begin{itemize}
        \item Analyze CETC's existing training/data pipeline
        \item Create a new GUI to interface with CETC's code structure
        \item Ensure compatibility with CETC's dataset format and workflow
        \item Resolve any dependency or conflicts
    \end{itemize}

    \item \textbf{Signed Receipt System} (7 hrs)
    \begin{itemize}
        \item Implement symmetric encryption with a moderator-held secret key
        \item Receipt captures: hyperparameters, class names, final mAP per class, model file size, timestamp, and unique run ID
        \item Receipt exported as encrypted \texttt{.receipt} file alongside model weights
        \item Moderator-side decryption utility for verification
    \end{itemize}

    \item \textbf{mAP and Model-Size Extraction} (3 hrs)
    \begin{itemize}
        \item Programmatically extract mAP from training results CSV
        \item Compute and record model file size (bytes)
        \item Feed both values into the receipt and display in GUI summary
    \end{itemize}

    \item \textbf{UI Adjustments and Testing} (6 hrs)
    \begin{itemize}
        \item Update dataset selection flow for CETC's folder structure
        \item Add receipt export status indicator in GUI
        \item End-to-end testing: train $\rightarrow$ export $\rightarrow$ receipt generation $\rightarrow$ decryption verification
    \end{itemize}
    \item \textbf{Web application deployment} (6 hrs)
    \begin{itemize}
        \item Deployment of the finalized, tested web application using the server with 1/8 Nvidia A10 GPU
    \end{itemize}
\end{enumerate}

\subsection*{Hour Breakdown}

\begin{center}
\begin{tabular}{lcc}
\toprule
\textbf{Task} & \textbf{Hours (est.)} \\
\midrule
Planning \& documentation & 5.0  \\
CETC code analysis \& integration & 15.0  \\
Signed receipt system (Fernet) & 7.0  \\
mAP / model-size extraction & 3.0  \\
UI adjustments \& testing & 6.0  \\
Web application deployment & 6.0  \\
\midrule
\textbf{Phase 2 Total} & \textbf{42.0}  \\
\bottomrule
\end{tabular}
\end{center}

%===============================================================================
\section*{Phase 3: XRP Robot Code and Testing}
%===============================================================================

\noindent \textbf{Estimated hours:} 23

\vspace{0.5em}

Testing and adaptation of the final part of the pipeline: utilizing the XRP robot with the model, the 3 object classes in a competition-ready environment. 

\subsection*{Scope of Work}

\begin{enumerate}[leftmargin=*]
    \item \textbf{State Machine Adaptation} (6 hrs)
    \begin{itemize}
        \item Align search/center/contact/avoid behavior with defined field dimensions
        \item Parameterize hardcoded values (distances, speeds, delays)
        \item Handle edge cases
    \end{itemize}

    \item \textbf{Parameter Tuning and Testing} (5 hrs)
    \begin{itemize}
        \item Tune turn speeds, centering threshold, approach distances for physical field
        \item Test with actual objects (cans, wheels, rocket model)
        \item Validate rangefinder accuracy at competition distances
    \end{itemize}

    \item \textbf{Integration Testing} (12 hrs)
    \begin{itemize}
        \item Full pipeline test: GUI-trained model $\rightarrow$ \texttt{run\_model.py} $\rightarrow$ XRP
        \item Verify ACK flow control stability over full 120-second run
        \item Confirm scoring conditions are met
    \end{itemize}
\end{enumerate}

\subsection*{Hour Breakdown}

\begin{center}
\begin{tabular}{lcc}
\toprule
\textbf{Task} & \textbf{Hours (est.)} \\
\midrule
State machine adaptation & 6.0  \\
Parameter tuning \& testing & 5.0  \\
Integration testing (full pipeline) & 12.0  \\
\midrule
\textbf{Phase 3 Total} & \textbf{23.0} \\
\bottomrule
\end{tabular}
\end{center}

%===============================================================================
\section*{Total Summary}
%===============================================================================

\begin{center}
\begin{tabular}{lccc}
\toprule
\textbf{Phase} & \textbf{Hours} \\
\midrule
2. GUI Creation & 42.0   \\
3. XRP Robot Code & 23.0   \\
\midrule
\textbf{Total} & \textbf{65.0}  \\
\bottomrule
\end{tabular}
\end{center}



\end{document}
